\section{Busca Tabu}
\label{sec:tabu}

A Busca Tabu é uma meta-heurística baseada em busca local com memória de curto prazo, projetada para evitar ciclos curtos e escapar de ótimos locais \cite{glover1989tabu,glover1990tabu,glover1997tabu}.
No VRPTW, versões clássicas utilizam um conjunto de movimentos factíveis, lista tabu por atributos e critério de aspiração.

\subsection{Estrutura implementada}

A versão final implementada segue uma linha clássica enxuta:
\begin{itemize}
  \item inicialização em solução I1;
  \item varredura completa de vizinhanças por iteração;
  \item memória tabu por cliente movimentado;
  \item aspiração global por melhoria do melhor incumbente;
  \item parada por limite de iterações ou estagnação.
\end{itemize}

As vizinhanças avaliadas em cada iteração são:
\begin{enumerate}
  \item \textbf{relocate} (intra e inter-rotas);
  \item \textbf{swap\_intra};
  \item \textbf{or-opt} (blocos de tamanho 2 ou 3, intra e inter-rotas).
\end{enumerate}

\begin{algorithm}[htbp]
\caption{Busca Tabu clássica (implementação final)}
\label{alg:tabu}
\begin{algorithmic}[1]
\Require instância VRPTW, tenure fixo $\theta$ ou intervalo $[\theta_{\min},\theta_{\max}]$, limites $I_{\max}$ e $L$
\Ensure melhor solução $S^*$
\State $S \leftarrow I1$, $S^* \leftarrow S$
\State inicializa vetor $tabuUntil[c] \leftarrow 0$
\State $noImprove \leftarrow 0$
\For{$t=1$ até $I_{\max}$}
  \State gera candidatos factíveis em relocate, swap\_intra e or-opt
  \State filtra candidatos tabu (exceto aspiração)
  \State escolhe melhor movimento admissível por custo (desempate por rotas)
  \If{não existe movimento admissível}
    \State \textbf{break}
  \EndIf
  \State aplica movimento em $S$
  \State amostra $\theta_t \sim U[\theta_{\min},\theta_{\max}]$ (ou $\theta_t \leftarrow \theta$ se fixo)
  \State para cada cliente atributo do movimento: $tabuUntil[c] \leftarrow t + \theta_t$
  \If{$S$ melhora $S^*$}
    \State $S^* \leftarrow S$; $noImprove \leftarrow 0$
  \Else
    \State $noImprove \leftarrow noImprove + 1$
  \EndIf
  \If{$noImprove \ge L$}
    \State \textbf{break}
  \EndIf
\EndFor
\State \Return $S^*$
\end{algorithmic}
\end{algorithm}

\paragraph{Texto complementar ao pseudocódigo.}
As linhas 5--7 explicitam a principal diferença para um método de descida: a Tabu escolhe o melhor movimento admissível da iteração, mesmo quando esse movimento não melhora a solução corrente.
Esse mecanismo controlado de piora temporária é o que permite escapar de ótimos locais.

Na implementação, cada candidato é representado por um descritor de movimento (tipo, rotas afetadas e clientes-atributo), solução resultante e avaliação lexicográfica (custo e número de rotas).
A linha 9 atualiza o bloqueio tabu no vetor \texttt{tabuUntil} para os clientes-atributo do movimento aplicado.
Com isso, evita-se reversão imediata do movimento recém-feito, reduzindo ciclos curtos.

As linhas 10--15 implementam o monitor de progresso: melhora do incumbente zera estagnação; ausência de melhora incrementa \texttt{noImprove}; o laço encerra quando atinge $L$.
Essa regra torna o custo computacional mais previsível sem descaracterizar a versão clássica.

\subsection{Regra tabu e aspiração}

Um movimento é considerado tabu se algum cliente atributo ainda estiver bloqueado:
\[
 tabuUntil[c] > t.
\]

O critério de aspiração libera movimento tabu quando ele melhora a melhor solução global $S^*$.
Esse desenho é compatível com a formulação clássica e evita bloqueios excessivos em movimentos realmente promissores.
Além disso, quando todos os movimentos da iteração estão tabu e nenhum satisfaz a aspiração global, aplica-se \emph{aspiration by default}: seleciona-se o melhor movimento tabu disponível para evitar estagnação por bloqueio total.

Nesta implementação comparativa, o tenure é mantido fixo por decisão de isonomia experimental entre métodos e para reduzir variância estocástica entre execuções.

\subsection{Parâmetros}

Os parâmetros calibráveis desta versão são:
\begin{itemize}
  \item $\theta$ (tenure fixo);
  \item $I_{\max}$ (máximo de iterações);
  \item $L$ (máximo de iterações sem melhoria).
\end{itemize}

Não foram incluídas extensões de memória de longo prazo (frequência, intensificação/diversificação), oscilação estratégica ou função de mérito penalizada.
A decisão metodológica foi manter uma versão clássica de curto prazo, comparável em pé de igualdade com VND e GRASP no estudo principal.

\subsection{Relação com o objetivo comparativo do trabalho}

Como a busca inicia em I1 (mesmo baseline dos outros métodos), o desempenho da Tabu observado nos experimentos representa ganho efetivo da estratégia de exploração com memória, e não vantagem de inicialização.
Essa separação é essencial para a análise comparativa entre VND, GRASP e Tabu.
